% GENERATED FILE. Edit canonical lessons, then run npm run generate:books.
% canonical-source-hash: fnv1a64:08c83e2b37b5ea57
% canonical-lessons: GU-C460-bholu, GU-C460-luccu, GU-C460-calak, GU-C460-kapat, GU-C460-phoslavvu

\chapter{Innocent or Sly}
\label{ch:gu-innocent-or-sly}
\hlchaptermodality{460}

\begin{chapteropening}
\textbf{By the end of this chapter:} \emph{I can say whether someone is easily fooled.}

Complete the last lesson of chapter 460: I can say whether someone is easily fooled.
\end{chapteropening}

\section[\texorpdfstring{\gu{ભોળું} (bhoḷũ)}{ભોળું (bhoḷũ)}]{\gu{ભોળું} (bhoḷũ) — innocent; trusting}

\label{lesson:GU-C460-bholu}



\begin{quote}
Before the new one: say the Gujarati for good-hearted, then the Gujarati for to soften with pity.
\end{quote}

\subsection*{You'll want to know: \gu{ભોળું}}
\textbf{\gu{ભોળું}} — \emph{bhoḷũ} — \textquotedblleft{}innocent; trusting\textquotedblright{}.

\begin{quote}
\gu{નાનું} \gu{બાળક} \gu{ભોળું} \gu{હોય} \gu{છે}.
\end{quote}

\begin{quote}
\emph{nānũ bāḷak bhoḷũ hoy che.}
\end{quote}

\begin{quote}
A small child is trusting.
\end{quote}

\begin{cousinweb}
\textbf{\gu{ભોળું}} (\emph{bhoḷũ}) goes back to a reconstructed \emph{*bhulla-}, \textquotedblleft{}defective\textquotedblright{}; CDIAL lists the Gujarati word, \textquotedblleft{}foolish, credulous\textquotedblright{} (CDIAL 9539), and today it is mostly the innocence.
\end{cousinweb}

\begin{grammarlens}[title={What you've built}]
One more word for describing what you lived and felt.
\end{grammarlens}

\subsection*{Your turn}
\begin{itemize}
\raggedright
  \item \emph{Say it:} \emph{bhoḷũ}
  \item \emph{Say it:} \emph{bhoḷũ}, once more
  \item \emph{Say it:} \emph{nānũ bāḷak bhoḷũ hoy che.} — A small child is trusting.
  \item \emph{Recall:} say the Gujarati for good-hearted, then the Gujarati for to soften with pity, then say \emph{bhoḷũ} again
\end{itemize}

\begin{tcolorbox}[breakable,colback=teal!4,colframe=teal!35!black,title={Before you move on}]
What does \gu{ભોળું} mean? (\textquotedblleft{}Innocent; trusting\textquotedblright{}.) Say it once more.
\end{tcolorbox}
\section[\texorpdfstring{\gu{લુચ્ચું} (luccũ)}{લુચ્ચું (luccũ)}]{\gu{લુચ્ચું} (luccũ) — sly; crafty}

\label{lesson:GU-C460-luccu}



\begin{quote}
Before the new one: say the Gujarati for to soften with pity, then the Gujarati for trusting.
\end{quote}

\subsection*{You'll want to know: \gu{લુચ્ચું}}
\textbf{\gu{લુચ્ચું}} — \emph{luccũ} — \textquotedblleft{}sly; crafty\textquotedblright{}.

\begin{quote}
\gu{તે} \gu{લુચ્ચો} \gu{વેપારી} \gu{બધાને} \gu{છેતરે} \gu{છે}.
\end{quote}

\begin{quote}
\emph{te lucco vepārī badhāne chetre che.}
\end{quote}

\begin{quote}
That sly trader cheats everyone.
\end{quote}

\begin{cousinweb}
\textbf{\gu{લુચ્ચું}} (\emph{luccũ}) goes back to a reconstructed \emph{*lucca-}, \textquotedblleft{}defective\textquotedblright{}; CDIAL lists the Gujarati word with the gloss \textquotedblleft{}lewd\textquotedblright{} (CDIAL 11073). Gujarati today uses it for someone sly, and no dictionary we use says how the two senses are linked.
\end{cousinweb}

\begin{grammarlens}[title={What you've built}]
One more word for describing what you lived and felt.
\end{grammarlens}

\subsection*{Your turn}
\begin{itemize}
\raggedright
  \item \emph{Say it:} \emph{luccũ}
  \item \emph{Say it:} \emph{luccũ}, once more
  \item \emph{Say it:} \emph{te lucco vepārī badhāne chetre che.} — That sly trader cheats everyone.
  \item \emph{Recall:} say the Gujarati for to soften with pity, then the Gujarati for trusting, then say \emph{luccũ} again
\end{itemize}

\begin{tcolorbox}[breakable,colback=teal!4,colframe=teal!35!black,title={Before you move on}]
What does \gu{લુચ્ચું} mean? (\textquotedblleft{}Sly; crafty\textquotedblright{}.) Say it once more.
\end{tcolorbox}
\section[\texorpdfstring{\gu{ચાલાક} (cālāk)}{ચાલાક (cālāk)}]{\gu{ચાલાક} (cālāk) — clever; shrewd}

\label{lesson:GU-C460-calak}



\begin{quote}
Before the new one: say the Gujarati for trusting, then the Gujarati for sly.
\end{quote}

\subsection*{You'll want to know: \gu{ચાલાક}}
\textbf{\gu{ચાલાક}} — \emph{cālāk} — \textquotedblleft{}clever; shrewd\textquotedblright{}.

\begin{quote}
\gu{શિયાળ} \gu{ઘણું} \gu{ચાલાક} \gu{પ્રાણી} \gu{છે}.
\end{quote}

\begin{quote}
\emph{śiyāḷ ghaṇũ cālāk prāṇī che.}
\end{quote}

\begin{quote}
The jackal is a very shrewd animal.
\end{quote}

\begin{cousinweb}
\textbf{\gu{ચાલાક}} (\emph{cālāk}) is usually traced to Persian \emph{čālāk}, \textquotedblleft{}quick, nimble, clever\textquotedblright{}.
\end{cousinweb}

\begin{grammarlens}[title={What you've built}]
One more word for describing what you lived and felt.
\end{grammarlens}

\subsection*{Your turn}
\begin{itemize}
\raggedright
  \item \emph{Say it:} \emph{cālāk}
  \item \emph{Say it:} \emph{cālāk}, once more
  \item \emph{Say it:} \emph{śiyāḷ ghaṇũ cālāk prāṇī che.} — The jackal is a very shrewd animal.
  \item \emph{Recall:} say the Gujarati for trusting, then the Gujarati for sly, then say \emph{cālāk} again
\end{itemize}

\begin{tcolorbox}[breakable,colback=teal!4,colframe=teal!35!black,title={Before you move on}]
What does \gu{ચાલાક} mean? (\textquotedblleft{}Clever; shrewd\textquotedblright{}.) Say it once more.
\end{tcolorbox}
\section[\texorpdfstring{\gu{કપટ} (kapaṭ)}{કપટ (kapaṭ)}]{\gu{કપટ} (kapaṭ) — deceit}

\label{lesson:GU-C460-kapat}



\begin{quote}
Before the new one: say the Gujarati for sly, then the Gujarati for shrewd.
\end{quote}

\subsection*{You'll want to know: \gu{કપટ}}
\textbf{\gu{કપટ}} — \emph{kapaṭ} — \textquotedblleft{}deceit\textquotedblright{}.

\begin{quote}
\gu{તેના} \gu{વચનમાં} \gu{કપટ} \gu{હતું}.
\end{quote}

\begin{quote}
\emph{tenā vacanmā̃ kapaṭ hatũ.}
\end{quote}

\begin{quote}
There was deceit in his promise.
\end{quote}

A neuter noun.

\begin{cousinweb}
\textbf{\gu{કપટ}} (\emph{kapaṭ}) is a tatsama, Sanskrit \emph{kapaṭa}, \textquotedblleft{}fraud, deceit\textquotedblright{}, as Monier-Williams records it.
\end{cousinweb}

\begin{grammarlens}[title={What you've built}]
One more word for describing what you lived and felt.
\end{grammarlens}

\subsection*{Your turn}
\begin{itemize}
\raggedright
  \item \emph{Say it:} \emph{kapaṭ}
  \item \emph{Say it:} \emph{kapaṭ}, once more
  \item \emph{Say it:} \emph{tenā vacanmā̃ kapaṭ hatũ.} — There was deceit in his promise.
  \item \emph{Recall:} say the Gujarati for sly, then the Gujarati for shrewd, then say \emph{kapaṭ} again
\end{itemize}

\begin{tcolorbox}[breakable,colback=teal!4,colframe=teal!35!black,title={Before you move on}]
What does \gu{કપટ} mean? (\textquotedblleft{}Deceit\textquotedblright{}.) Say it once more.
\end{tcolorbox}
\section[\texorpdfstring{\gu{ફોસલાવવું} (phoslāvvũ)}{ફોસલાવવું (phoslāvvũ)}]{\gu{ફોસલાવવું} (phoslāvvũ) — to coax; to cajole}

\label{lesson:GU-C460-phoslavvu}



\begin{quote}
Before the new one: say the Gujarati for shrewd, then the Gujarati for deceit.
\end{quote}

\subsection*{You'll want to know: \gu{ફોસલાવવું}}
\textbf{\gu{ફોસલાવવું}} — \emph{phoslāvvũ} — \textquotedblleft{}to coax; to cajole\textquotedblright{}.

\begin{quote}
\gu{માતા} \gu{બાળકને} \gu{મીઠાઈથી} \gu{ફોસલાવે} \gu{છે}.
\end{quote}

\begin{quote}
\emph{mātā bāḷakne mīṭhāīthī phoslāve che.}
\end{quote}

\begin{quote}
The mother coaxes the child with sweets.
\end{quote}

Past, with \gu{તેણે} (he or she …): \gu{તેણે} \gu{બાળકને} \gu{ફોસલાવ્યું}, \emph{teṇe bāḷakne phoslāvyũ}.

\begin{cousinweb}
\textbf{\gu{ફોસલાવવું}} (\emph{phoslāvvũ}) is inherited: CDIAL lists the Gujarati verb, \textquotedblleft{}to cajole\textquotedblright{}, under Sanskrit \emph{spṛśati}, \textquotedblleft{}touches, strokes\textquotedblright{} (CDIAL 13815).
\end{cousinweb}

\begin{grammarlens}[title={What you've built}]
One more word for describing what you lived and felt.
\end{grammarlens}

\subsection*{Your turn}
\begin{itemize}
\raggedright
  \item \emph{Say it:} \emph{phoslāvvũ}
  \item \emph{Say it:} \emph{phoslāvvũ}, once more
  \item \emph{Say it:} \emph{mātā bāḷakne mīṭhāīthī phoslāve che.} — The mother coaxes the child with sweets.
  \item \emph{Recall:} say the Gujarati for shrewd, then the Gujarati for deceit, then say \emph{phoslāvvũ} again
\end{itemize}

\begin{tcolorbox}[breakable,colback=teal!4,colframe=teal!35!black,title={Before you move on}]
What does \gu{ફોસલાવવું} mean? (\textquotedblleft{}To coax; to cajole\textquotedblright{}.) Say it once more.
\end{tcolorbox}
